\section{Tổng quan về mã hóa trong ứng dụng nhắn tin}

\subsection{Mô hình đe dọa trong giao tiếp mạng}

Khi một tin nhắn được gửi đi giữa hai người dùng, dữ liệu phải đi qua nhiều thành phần trung gian trước khi đến được đích. Trên đường truyền từ thiết bị của người gửi đến máy chủ của nhà cung cấp dịch vụ, gói tin có thể bị các điểm trung gian như bộ định tuyến, nhà cung cấp dịch vụ mạng hoặc các thiết bị phân tích lưu lượng theo dõi. Ngay cả khi dữ liệu đến được máy chủ, nó vẫn có nguy cơ bị chính nhà cung cấp dịch vụ đọc, lưu trữ hoặc chuyển giao cho bên thứ ba. Đây là thực tế đã được ghi nhận trong nhiều vụ việc lộ dữ liệu người dùng từ các nền tảng nhắn tin phổ biến trong những năm gần đây.

\textbf{Tài sản cần bảo vệ.} Trong một ứng dụng nhắn tin, các tài sản chính gồm: (i) nội dung tin nhắn và tệp đính kèm; (ii) metadata như ai liên lạc với ai, vào thời điểm nào và với tần suất ra sao; (iii) danh tính và vật liệu khóa của người dùng gồm khóa định danh, các prekey và trạng thái phiên; và (iv) thông tin xác thực tài khoản như mật khẩu. Ưu tiên cao nhất là bảo vệ nội dung và vật liệu khóa, vì để lộ chúng đồng nghĩa với mất hoàn toàn tính bí mật của hội thoại.

\textbf{Các loại đối thủ.} Mô hình đe dọa chính của hệ thống nhắn tin xét các đối thủ sau:
\begin{itemize}
  \item Đối thủ nghe lén thụ động trên đường truyền (nhà cung cấp dịch vụ mạng, thiết bị trung gian, hay một chính phủ có quyền giám sát hạ tầng): đọc và lưu lại mọi gói tin nhưng không chủ động sửa đổi.
  \item Đối thủ chủ động thực hiện tấn công xen giữa: chặn, sửa, chèn hoặc phát lại gói tin, và có thể mạo danh một đầu nếu cơ chế xác thực yếu.
  \item Máy chủ vẫn tuân thủ giao thức để duy trì dịch vụ nhưng cố suy luận nội dung từ dữ liệu mà nó lưu trữ và trung chuyển. Đây chính là đối thủ trung tâm mà mã hóa đầu cuối nhắm tới, vì máy chủ có khả năng truy cập toàn bộ dữ liệu đi qua nó nếu không áp dụng mã hóa đầu cuối.
  \item Đối thủ tương lai có máy tính lượng tử thu thập bản mã hôm nay theo mô hình ``harvest now, decrypt later'' và chờ giải mã khi có đủ năng lực lượng tử.
\end{itemize}

\textbf{Giả định tin cậy.} Thiết bị đầu cuối của người dùng cùng các thư viện mật mã chạy trên đó được coi là đáng tin trong phạm vi đồ án; nếu chính thiết bị bị chiếm quyền thì khóa bí mật và bản rõ đều có thể bị lộ, điều này nằm ngoài khả năng bảo vệ của giao thức. Máy chủ không được tin để biết nội dung, nhưng vẫn được dựa vào cho tính sẵn sàng, thứ tự và việc trung chuyển tin; một máy chủ độc hại có thể chặn hoặc làm chậm tin (tấn công vào tính sẵn sàng) nhưng vẫn không đọc được nội dung đã mã hóa đầu cuối.

\subsection{Các cấp độ bảo vệ dữ liệu}

Trong thực tế, các ứng dụng nhắn tin hiện nay áp dụng nhiều cấp độ bảo vệ khác nhau với mức độ hiệu quả không đồng đều. Cấp độ cơ bản nhất là mã hóa tầng vận chuyển, thường được hiện thực hóa qua TLS. Cơ chế này chỉ bảo vệ dữ liệu trên đường truyền từ thiết bị tới máy chủ, đồng nghĩa với việc máy chủ vẫn có khả năng đọc được toàn bộ nội dung tin nhắn trước khi chuyển tiếp.

Cấp độ cao hơn là mã hóa phía máy chủ, trong đó dữ liệu được lưu trữ dưới dạng mã hóa nhưng chính máy chủ nắm giữ khóa giải mã. Phương án này bảo vệ tốt trước các rủi ro về rò rỉ dữ liệu do xâm phạm phần cứng hoặc bản sao lưu, nhưng vẫn không ngăn được việc nhân viên nội bộ hoặc cơ quan chức năng truy cập nội dung khi có yêu cầu.

Cấp độ bảo vệ cao nhất và cũng là tiêu chuẩn trong các hệ thống nhắn tin hiện đại chính là mã hóa đầu cuối. Trong mô hình này, tin nhắn được mã hóa ngay trên thiết bị của người gửi bằng khóa mà chỉ người nhận có thể giải mã. Máy chủ chỉ đóng vai trò trung chuyển các bản mã mà không có bất kỳ khả năng đọc nội dung nào. Đây là mô hình được Signal tiên phong và sau đó được WhatsApp, iMessage và nhiều nền tảng khác áp dụng.

\subsection{Yêu cầu đối với một hệ thống mã hóa đầu cuối}

Một hệ thống mã hóa đầu cuối an toàn cần đáp ứng đồng thời nhiều thuộc tính an toàn. Tính bí mật đảm bảo rằng chỉ người nhận hợp lệ mới có thể đọc được nội dung tin nhắn. Tính toàn vẹn bảo đảm rằng nội dung không bị sửa đổi trong quá trình truyền tải, đồng thời cho phép người nhận phát hiện mọi thay đổi dù nhỏ nhất. Tính xác thực cho phép người nhận xác minh chắc chắn rằng tin nhắn xuất phát từ người gửi được tuyên bố. Ngoài ra, các hệ thống tiên tiến còn hướng tới thuộc tính forward secrecy, nghĩa là ngay cả khi khóa dài hạn bị lộ trong tương lai, các phiên tin nhắn trong quá khứ vẫn không thể bị giải mã.

Để hiện thực hóa các yêu cầu trên, hệ thống cần kết hợp giữa mã hóa đối xứng cho nội dung tin nhắn với một cơ chế trao đổi khóa an toàn giữa hai đầu. Chính cơ chế trao đổi khóa là nơi mà các thuật toán mật mã khóa công khai truyền thống đang đối mặt với nguy cơ trực tiếp từ máy tính lượng tử, đòi hỏi phải được thay thế bằng những giải pháp kháng lượng tử.

\section{Mật mã đối xứng và mã hóa có xác thực}

\subsection{Thuật toán AES và chế độ GCM}

AES là một thuật toán mã hóa khối được Viện Tiêu chuẩn và Công nghệ Quốc gia Hoa Kỳ chuẩn hóa từ năm 2001 và hiện vẫn là thuật toán mã hóa đối xứng phổ biến nhất trên thế giới. AES hoạt động trên các khối dữ liệu 128-bit với độ dài khóa 128, 192 hoặc 256 bit, trong đó biến thể AES-256 cung cấp mức an toàn cao nhất. Thuật toán bao gồm nhiều vòng lặp với các thao tác thay thế, dịch hàng, trộn cột và cộng khóa vòng, tạo ra một hoán vị phụ thuộc khóa có tính kháng tấn công cao.

AES ở dạng nguyên bản chỉ định nghĩa phép mã hóa trên một khối 128-bit, do đó trong thực tế cần kết hợp với một chế độ vận hành để xử lý các thông điệp có độ dài bất kỳ. Trong các ứng dụng hiện đại, chế độ GCM là lựa chọn được khuyến nghị bởi NIST nhờ khả năng cung cấp đồng thời mã hóa và xác thực trong một thao tác duy nhất. Cấu trúc này được gọi chung là AEAD, viết tắt của Authenticated Encryption with Associated Data.

Khi vận hành AES-GCM, ngoài khóa và bản rõ, thuật toán còn cần một giá trị nonce 96-bit. Nonce phải là duy nhất đối với mỗi lần mã hóa dưới cùng một khóa, nhưng không cần bí mật và có thể được truyền công khai cùng với bản mã. Đầu ra của thuật toán bao gồm bản mã có độ dài bằng bản rõ và một thẻ xác thực 128-bit. Phía giải mã sẽ kiểm tra thẻ này trước khi trả về bản rõ, bất kỳ sự thay đổi nào trong bản mã hoặc nonce đều dẫn đến việc thẻ không hợp lệ và quá trình giải mã sẽ thất bại. Đây là một tính chất rất quan trọng giúp hệ thống tự động phát hiện tấn công sửa đổi dữ liệu mà không cần cơ chế xác thực riêng.

\subsection{Hàm dẫn xuất khóa từ mật khẩu}

Trong các ứng dụng nhắn tin sử dụng tài khoản, mật khẩu là thông tin duy nhất người dùng có thể ghi nhớ và sử dụng để thiết lập truy cập. Tuy nhiên, mật khẩu thường có entropy thấp và không thể được sử dụng trực tiếp làm khóa mã hóa, vì vậy cần một hàm dẫn xuất khóa để chuyển đổi mật khẩu thành khóa có kích thước và entropy phù hợp.

PBKDF2 là hàm dẫn xuất khóa được chuẩn hóa trong RFC 8018 và được sử dụng rộng rãi trong nhiều hệ thống bảo mật. PBKDF2 nhận đầu vào gồm mật khẩu, một giá trị salt, số vòng lặp và hàm băm cơ sở, sau đó trả về khóa dẫn xuất có độ dài mong muốn. Salt đóng vai trò quan trọng trong việc ngăn chặn tấn công sử dụng bảng băm tiền tính toán, còn số vòng lặp cao làm tăng chi phí tính toán của mỗi lần thử mật khẩu, khiến các tấn công brute-force trở nên tốn kém. Hướng dẫn lưu trữ mật khẩu của OWASP đề xuất sử dụng ít nhất 600.000 vòng khi kết hợp PBKDF2 với HMAC-SHA-256, trong khi hướng dẫn về định danh số của NIST (SP 800-63B) đưa ra các yêu cầu chung về độ mạnh và chính sách mật khẩu; cả hai cùng nhằm cân bằng giữa an toàn và trải nghiệm người dùng.

Một đặc điểm quan trọng của PBKDF2 là tính xác định. Khi đầu vào gồm mật khẩu, salt, số vòng và hàm băm không đổi, đầu ra luôn là cùng một khóa. Tính chất này cho phép ứng dụng dẫn xuất lại khóa chủ mỗi khi người dùng đăng nhập, mà không cần lưu khóa này ở bất cứ đâu.

\section{Mật mã khóa công khai truyền thống và mối đe dọa từ máy tính lượng tử}

\subsection{Nguyên lý mã hóa khóa công khai}

Khác với mã hóa đối xứng trong đó cả hai bên dùng chung một khóa bí mật, mật mã khóa công khai sử dụng một cặp khóa gồm khóa công khai và khóa bí mật. Khóa công khai có thể được phát hành rộng rãi và cho phép bất kỳ ai mã hóa thông điệp hoặc xác minh chữ ký số, trong khi khóa bí mật chỉ được người sở hữu giữ lại và dùng để giải mã hoặc ký. Cơ chế này giải quyết được bài toán trao đổi khóa trong môi trường mở, vốn là trở ngại lớn nhất của mã hóa đối xứng.

An toàn của các hệ mã hóa khóa công khai dựa trên các bài toán được cho là khó với máy tính cổ điển. Hệ RSA dựa trên độ khó của bài toán phân tích một số nguyên lớn thành tích các số nguyên tố. Các hệ mã đường cong elliptic như ECDH và ECDSA dựa trên bài toán logarit rời rạc trên nhóm điểm của đường cong elliptic, với độ khó tương đương một số bit gấp khoảng năm lần so với RSA ở cùng mức an toàn, cho phép khóa nhỏ gọn hơn đáng kể.

\subsection{Thuật toán Shor và hệ quả}

An toàn của RSA và ECDH dựa trên một giả thiết: các bài toán phân tích số nguyên và logarit rời rạc là khó đối với máy tính cổ điển. Thật vậy, cho đến nay chưa có thuật toán cổ điển nào giải được hai bài toán này trong thời gian đa thức, nên với kích thước khóa đủ lớn thì việc phá khóa bằng máy tính thông thường là bất khả thi trên thực tế. Năm 1994, Peter Shor công bố một thuật toán chạy trên máy tính lượng tử giải được cả hai bài toán đó trong thời gian đa thức, tức là phá vỡ chính giả thiết nền tảng nói trên. Hệ quả là một máy tính lượng tử đủ lớn và đủ ổn định để chạy được thuật toán Shor sẽ bẻ được khóa RSA và ECDH trong thời gian khả thi, làm sụp đổ nền tảng của toàn bộ mật mã khóa công khai đang được sử dụng hiện nay. Hệ quả này không chỉ dừng ở khâu trao đổi khóa: các lược đồ chữ ký số dựa trên cùng bài toán logarit rời rạc trên đường cong elliptic - như ECDSA hay EdDSA (Ed25519 là một biến thể phổ biến) - cũng bị Shor phá tương tự, nên cả việc xác thực bằng chữ ký số cũng không còn an toàn trước máy tính lượng tử.

Mặc dù các máy tính lượng tử hiện tại chưa đủ quy mô để thực thi Shor trên khóa có độ dài thực tế, xu hướng phát triển phần cứng lượng tử đang diễn ra rất nhanh, với các thông báo mới về số qubit và độ ổn định xuất hiện liên tục trong những năm gần đây. Chính vì vậy, cộng đồng nghiên cứu đã sớm đặt ra yêu cầu phải chuẩn bị các thuật toán thay thế trước khi máy tính lượng tử trở thành hiện thực.

Ngoài ra, có một lớp tấn công cần được tính đến ngay lập tức dù máy tính lượng tử chưa tồn tại, đó là mô hình tấn công “harvest now, decrypt later”. Trong mô hình này, kẻ tấn công thu thập và lưu trữ dữ liệu mã hóa hiện tại, chờ đến khi máy tính lượng tử đủ mạnh xuất hiện để tiến hành giải mã. Những dữ liệu có giá trị lâu dài như hồ sơ y tế, thông tin tình báo hay tài liệu pháp lý đặc biệt dễ trở thành mục tiêu. Đó là lý do các tiêu chuẩn mật mã hậu lượng tử cần được triển khai càng sớm càng tốt, ngay trong giai đoạn chuyển tiếp.

\section{Mật mã hậu lượng tử}


\subsection{Các họ thuật toán kháng lượng tử}

Mật mã hậu lượng tử là tập hợp các thuật toán được thiết kế để có thể chạy trên phần cứng cổ điển hiện có nhưng vẫn an toàn trước các tấn công từ cả máy tính cổ điển và máy tính lượng tử. Các thuật toán này được xây dựng trên các bài toán khó mà chưa có thuật toán lượng tử nào được biết đến có khả năng giải trong thời gian đa thức.

Cộng đồng nghiên cứu đã đề xuất nhiều họ thuật toán khác nhau. Họ dựa trên lưới (lattice-based) dựa trên độ khó của các bài toán trên lưới như bài toán học kèm lỗi (Learning With Errors, LWE) hay bài toán tìm vector ngắn nhất (Shortest Vector Problem, SVP), nổi bật với tính linh hoạt và hiệu năng tốt. Họ dựa trên mã (code-based) dựa trên độ khó của giải mã mã tuyến tính ngẫu nhiên, có lịch sử lâu đời từ hệ McEliece năm 1978 nhưng có nhược điểm là khóa công khai rất lớn. Họ dựa trên hàm băm (hash-based) xây dựng chữ ký số trên nền các hàm băm đã được chứng minh an toàn, ví dụ như SPHINCS+, có ưu điểm là độ tin cậy về mặt lý thuyết rất cao nhưng chữ ký lại tương đối lớn. Họ đa biến (multivariate) và họ isogeny đang được nghiên cứu nhưng còn nhiều thách thức về mặt an toàn và hiệu năng.

\subsection{Cơ sở toán học của mật mã trên lưới}

Trước khi đi vào thuật toán ML-KEM, cần nhắc lại một số khái niệm nền tảng về lưới (lattice), vốn là đối tượng toán học trung tâm cho gần như toàn bộ các đề xuất hậu lượng tử được NIST đánh giá cao nhất.

\textbf{Định nghĩa lưới (lattice).} Cho $n$ vector $\mathbf{b}_1, \mathbf{b}_2, \ldots, \mathbf{b}_n$ độc lập tuyến tính trong $\mathbb{R}^m$. Lưới $\mathcal{L}$ sinh bởi các vector này được định nghĩa là tập hợp tất cả các tổ hợp tuyến tính với hệ số nguyên:
\[
\mathcal{L}(\mathbf{b}_1, \ldots, \mathbf{b}_n) = \left\{ \sum_{i=1}^{n} x_i \mathbf{b}_i \,\middle|\, x_i \in \mathbb{Z} \right\}.
\]
Tập $\{\mathbf{b}_1, \ldots, \mathbf{b}_n\}$ được gọi là một cơ sở của lưới, còn $n$ là bậc của lưới. Về mặt hình học, lưới là một mạng các điểm phân bố đều trong không gian, và cùng một lưới có thể được biểu diễn bằng vô số cơ sở khác nhau.

\textbf{Các bài toán khó trên lưới.} Hai bài toán nền tảng được sử dụng làm cơ sở cho an toàn của mật mã trên lưới gồm:
\begin{itemize}
    \item \textbf{Bài toán tìm vector ngắn nhất (Shortest Vector Problem, SVP)}: cho một cơ sở của lưới $\mathcal{L}$, tìm vector khác không $\mathbf{v} \in \mathcal{L}$ có độ dài Euclid $\|\mathbf{v}\|$ nhỏ nhất.
    \item \textbf{Bài toán tìm vector gần nhất (Closest Vector Problem, CVP)}: cho một cơ sở của lưới $\mathcal{L}$ và một điểm $\mathbf{t} \in \mathbb{R}^m$ không thuộc lưới, tìm vector $\mathbf{v} \in \mathcal{L}$ sao cho $\|\mathbf{v} - \mathbf{t}\|$ nhỏ nhất.
\end{itemize}
Cả hai bài toán đều đã được chứng minh là NP-hard trong trường hợp chính xác và vẫn rất khó ngay cả với các phiên bản xấp xỉ. Điều đặc biệt quan trọng là cho đến nay, không có thuật toán lượng tử nào được biết giải SVP hoặc CVP trong thời gian đa thức, khác hoàn toàn với bài toán phân tích số nguyên hay logarit rời rạc mà thuật toán Shor đã giải được.

\textbf{Bài toán học kèm lỗi (Learning With Errors, LWE).} Bài toán LWE do Oded Regev đề xuất năm 2005 là cầu nối giữa các bài toán trên lưới dạng hình học nêu trên và mật mã ứng dụng. Cho $n$, $q$ là các số nguyên dương và $\chi$ là một phân phối xác suất trên $\mathbb{Z}_q$ trả về các giá trị nhỏ. Với vector bí mật $\mathbf{s} \in \mathbb{Z}_q^n$, phân phối LWE sinh các mẫu có dạng
\[
(\mathbf{a}, b) = \left(\mathbf{a}, \, \langle \mathbf{a}, \mathbf{s} \rangle + e \bmod q \right),
\]
trong đó $\mathbf{a} \in \mathbb{Z}_q^n$ được chọn đều ngẫu nhiên và $e \leftarrow \chi$ là nhiễu nhỏ. Bài toán \emph{search-LWE} yêu cầu khôi phục $\mathbf{s}$ khi biết nhiều mẫu $(\mathbf{a}_i, b_i)$, còn bài toán \emph{decision-LWE} yêu cầu phân biệt phân phối các mẫu LWE với phân phối đều. Khi viết dưới dạng ma trận với $m$ mẫu, ta có
\[
\mathbf{b} = A \cdot \mathbf{s} + \mathbf{e} \bmod q,
\]
với $A \in \mathbb{Z}_q^{m \times n}$ ngẫu nhiên và $\mathbf{e} \in \mathbb{Z}_q^m$ là vector nhiễu có các phần tử nhỏ. Nhờ công trình của Regev, bài toán decision-LWE có mức độ khó được chứng minh bằng reduction tới SVP trong trường hợp xấu nhất, nghĩa là nếu LWE giải được trong thời gian trung bình thì SVP cũng giải được trong thời gian tồi nhất. Đây là tính chất rất mạnh mà rất ít giả thiết mật mã khác có được.

\textbf{Bài toán Module-LWE.} Nếu thay không gian làm việc từ $\mathbb{Z}_q$ sang vành đa thức $R_q = \mathbb{Z}_q[X] / (X^n + 1)$ với $n$ là một lũy thừa của $2$, ta thu được biến thể Module-LWE. Các phần tử của $R_q$ là các đa thức có bậc nhỏ hơn $n$ với hệ số trong $\mathbb{Z}_q$, và phép nhân trong $R_q$ được thực hiện modulo đa thức $X^n + 1$. Bài toán Module-LWE cho một ma trận $\mathbf{A} \in R_q^{k \times k}$ và vector $\mathbf{b} = \mathbf{A} \cdot \mathbf{s} + \mathbf{e} \in R_q^k$ yêu cầu khôi phục $\mathbf{s} \in R_q^k$. Về bản chất Module-LWE là LWE đóng gói lại theo khối: thay vì một số nguyên bí mật ta có một đa thức bí mật có $n$ hệ số, và các phép nhân được thực hiện trên đa thức chứ không phải số nguyên riêng lẻ.

Ưu điểm lớn của Module-LWE là giảm đáng kể kích thước khóa so với LWE thuần, vì một phần tử của $R_q$ đã mang thông tin tương đương $n$ phần tử của $\mathbb{Z}_q$, trong khi vẫn giữ được độ khó tương đương nhờ cấu trúc đại số phong phú của vành đa thức. Đồng thời, các phép nhân đa thức có thể được tăng tốc đáng kể bằng biến đổi NTT (Number Theoretic Transform), giúp Module-LWE đạt được hiệu năng gần với các hệ mã hóa đối xứng.

\subsection{Chương trình chuẩn hóa của NIST}

Nhận thức được tầm quan trọng của việc chuẩn hóa mật mã hậu lượng tử, Viện Tiêu chuẩn và Công nghệ Quốc gia Hoa Kỳ đã khởi động chương trình tuyển chọn thuật toán từ năm 2016. Chương trình trải qua ba vòng đánh giá công khai với sự tham gia của hàng trăm nhà nghiên cứu trên toàn thế giới, mỗi thuật toán đều được xem xét về mặt toán học, hiệu năng cũng như khả năng cài đặt an toàn.

Năm 2022, NIST thông báo kết thúc vòng tuyển chọn chính và công bố các thuật toán được chọn để chuẩn hóa. Trong đó, Kyber được chọn làm thuật toán trao đổi khóa chuẩn cho mục đích mã hóa thông tin công cộng, Dilithium được chọn cho chữ ký số tổng quát, Falcon cho các trường hợp cần chữ ký ngắn, và SPHINCS+ được duy trì như một phương án dự phòng dựa trên hàm băm. Đến năm 2024, NIST chính thức công bố các tiêu chuẩn FIPS 203, 204 và 205, trong đó Kyber được chuẩn hóa dưới tên gọi ML-KEM.

Để thống nhất cách gọi trong phần còn lại của đồ án: ``Kyber'' (đầy đủ là CRYSTALS-Kyber) là tên thuật toán gốc, còn ``ML-KEM'' là tên đã được chuẩn hóa trong FIPS 203, với ba mức tham số ML-KEM-512, ML-KEM-768 và ML-KEM-1024 (tương ứng Kyber-512/768/1024). Hai tên chỉ cùng một họ thuật toán nên báo cáo dùng tên chuẩn ML-KEM và chỉ nhắc tới Kyber khi nói về thiết kế gốc. Một tên thứ ba là X-Wing, một KEM hybrid kết hợp X25519 với ML-KEM-768 được đồ án dùng cho hội thoại nhóm; khác với ML-KEM đã được chuẩn hóa, X-Wing hiện vẫn là bản nháp của IETF (\texttt{draft-connolly-cfrg-xwing-kem}) mang tính thực nghiệm, chưa phải một chuẩn đã ổn định.

\subsection{Cơ chế Key Encapsulation}

Một điểm khác biệt quan trọng của ML-KEM và nhiều thuật toán hậu lượng tử khác so với RSA là chúng thường được thiết kế dưới dạng Key Encapsulation Mechanism thay vì mã hóa bất đối xứng trực tiếp. Một KEM là bộ ba thuật toán $(\text{KeyGen}, \text{Encaps}, \text{Decaps})$ trong đó:
\begin{itemize}
    \item $\text{KeyGen}() \to (pk, sk)$: sinh cặp khóa công khai và bí mật.
    \item $\text{Encaps}(pk) \to (c, K)$: người gửi cung cấp khóa công khai của người nhận, nhận về một ciphertext $c$ và một shared secret $K$.
    \item $\text{Decaps}(sk, c) \to K$: người nhận dùng khóa bí mật của mình và ciphertext vừa nhận để khôi phục cùng shared secret $K$.
\end{itemize}
Shared secret $K$ sau đó được dùng làm khóa cho một thuật toán đối xứng như AES-GCM để mã hóa dữ liệu thực tế. Ưu điểm của cấu trúc KEM là đơn giản, có chứng minh bảo mật chặt chẽ theo mô hình IND-CCA2 và phù hợp với đặc điểm của nhiều bài toán trên lưới, nơi mà việc mã hóa trực tiếp một thông điệp tùy ý thường gặp nhiều hạn chế kỹ thuật.

\section{Thuật toán ML-KEM (Kyber)}

\subsection{Không gian làm việc và các tham số}

ML-KEM hoạt động trên vành đa thức
\[
R_q = \mathbb{Z}_q[X] / (X^n + 1),
\]
với các tham số cố định $n = 256$ và $q = 3329$. Mỗi phần tử của $R_q$ là một đa thức dạng $a_0 + a_1 X + a_2 X^2 + \cdots + a_{255} X^{255}$ trong đó các hệ số $a_i \in \mathbb{Z}_q = \{0, 1, \ldots, 3328\}$. Phép cộng trong $R_q$ đơn giản là cộng hệ số theo modulo $q$, còn phép nhân được thực hiện bằng cách nhân đa thức thông thường rồi rút gọn modulo $X^n + 1$. Do $X^n \equiv -1 \pmod{X^n+1}$, phép rút gọn này chỉ đơn thuần là hoán vị dấu các hệ số bậc cao, khiến việc cài đặt rất hiệu quả.

Tham số $q = 3329$ là một số nguyên tố đặc biệt thỏa mãn $q \equiv 1 \pmod{2n}$, cho phép vành $R_q$ tồn tại căn bậc $2n$ của đơn vị. Nhờ đó, phép nhân đa thức có thể được tăng tốc bằng biến đổi NTT, đưa độ phức tạp từ $O(n^2)$ xuống $O(n \log n)$, tương tự như vai trò của biến đổi Fourier rời rạc trong xử lý tín hiệu.

Các phiên bản ML-KEM khác nhau chỉ thay đổi tham số $k$, là số chiều của module làm việc. Cụ thể ma trận công khai $\mathbf{A}$ thuộc $R_q^{k \times k}$, còn các vector bí mật $\mathbf{s}$, nhiễu $\mathbf{e}$ và vector mã hóa thuộc $R_q^k$. Ba biến thể được sử dụng là:
\begin{itemize}
    \item ML-KEM-512 với $k = 2$, tương ứng mức an toàn NIST Level 1 (tương đương AES-128).
    \item ML-KEM-768 với $k = 3$, tương ứng mức an toàn NIST Level 3 (tương đương AES-192).
    \item ML-KEM-1024 với $k = 4$, tương ứng mức an toàn NIST Level 5 (tương đương AES-256).
\end{itemize}
Trong đồ án này, ML-KEM-768 được dùng làm ví dụ toán học vì đây là biến thể cân bằng và dễ trình bày. 

Bảng sau tóm tắt đầu vào, đầu ra và vai trò của ba thao tác chính trong ML-KEM. Đây cũng là giao diện mà ứng dụng dùng khi gọi thư viện liboqs, thay vì tự thao tác trực tiếp với đa thức bên trong thuật toán.

\begin{table}[H]
\centering
\renewcommand{\arraystretch}{1.25}
\begin{tabularx}{\linewidth}{|p{3cm}|p{3.8cm}|p{3.8cm}|X|}
\hline
\textbf{Thao tác} & \textbf{Đầu vào} & \textbf{Đầu ra} & \textbf{Vai trò trong hệ thống} \\
\hline
\texttt{KeyGen} & Nguồn ngẫu nhiên an toàn và bộ tham số thuật toán & Khóa công khai $pk$ và khóa bí mật $sk$ & Người nhận công bố $pk$ trong prekey bundle, còn $sk$ chỉ lưu cục bộ trên client. \\
\hline
\texttt{Encaps} & Khóa công khai $pk$ của người nhận & Ciphertext KEM $c$ và shared secret $K$ & Người gửi tạo một secret chung mà chỉ người giữ $sk$ tương ứng mới khôi phục được. \\
\hline
\texttt{Decaps} & Khóa bí mật $sk$ và ciphertext KEM $c$ & Shared secret $K$ hoặc một secret giả ngẫu nhiên nếu ciphertext không hợp lệ & Người nhận khôi phục secret chung để đưa vào HKDF và thiết lập khóa phiên. \\
\hline
\end{tabularx}
\caption{Đầu vào, đầu ra và vai trò của các thao tác ML-KEM}
\end{table}

\subsection{Phân phối nhiễu tâm nhị thức}

Các vector bí mật và nhiễu trong ML-KEM được lấy mẫu từ phân phối tâm nhị thức $B_\eta$, với $\eta$ là tham số nhỏ. Một mẫu $x \leftarrow B_\eta$ được sinh bằng cách lấy $2\eta$ bit ngẫu nhiên độc lập $b_1, \ldots, b_\eta, b'_1, \ldots, b'_\eta$ và tính:
\[
x = \sum_{i=1}^{\eta} (b_i - b'_i).
\]
Giá trị $x$ nằm trong khoảng $[-\eta, \eta]$ với phân phối đối xứng quanh $0$. Với ML-KEM-768, tham số $\eta_1 = 2$ được dùng cho $\mathbf{s}$ và $\mathbf{e}$ trong KeyGen, còn $\eta_2 = 2$ được dùng cho nhiễu trong Encaps. Khi mở rộng phân phối từ một số nguyên sang một đa thức thuộc $R_q$, mỗi hệ số của đa thức được lấy mẫu độc lập từ $B_\eta$, cho nên một ``đa thức nhỏ'' trong $R_q$ là một đa thức mà toàn bộ $n = 256$ hệ số đều có giá trị tuyệt đối nhỏ.

Việc chọn phân phối tâm nhị thức thay vì phân phối Gaussian rời rạc xuất phát từ nhu cầu cài đặt: $B_\eta$ chỉ cần các phép xor và đếm bit, hoàn toàn tránh được các phép tính dấu phẩy động vốn nhạy cảm với tấn công kênh kề. Về mặt bảo mật, $B_\eta$ cho độ khó Module-LWE gần như tương đương với Gaussian có phương sai tương ứng.

\subsection{Sinh khóa}

Thuật toán ML-KEM.KeyGen không nhận thông điệp của người dùng. Đầu vào thực tế của nó là nguồn ngẫu nhiên an toàn và bộ tham số như ML-KEM-768 hoặc ML-KEM-1024. Đầu ra là cặp khóa $(pk, sk)$, trong đó $pk$ được gửi lên server để người khác bắt tay, còn $sk$ phải được giữ bí mật ở client.

Các bước cốt lõi của quá trình sinh khóa như sau:

\begin{enumerate}
    \item Sinh một seed $\rho \in \{0,1\}^{256}$ và một seed $\sigma \in \{0,1\}^{256}$ độc lập, ngẫu nhiên.
    \item Mở rộng $\rho$ thành ma trận ngẫu nhiên $\mathbf{A} \in R_q^{k \times k}$ bằng hàm $\text{XOF}(\rho, i, j)$ (cụ thể là SHAKE-128), đảm bảo rằng từ cùng $\rho$ ta luôn dựng lại được cùng $\mathbf{A}$.
    \item Lấy mẫu vector bí mật $\mathbf{s} \in R_q^k$ với mỗi thành phần đa thức có các hệ số theo phân phối $B_{\eta_1}$, sử dụng seed $\sigma$ kết hợp bộ đếm.
    \item Lấy mẫu vector nhiễu $\mathbf{e} \in R_q^k$ cũng theo $B_{\eta_1}$ với seed $\sigma$ và bộ đếm tiếp theo.
    \item Tính
    \[
    \mathbf{t} = \mathbf{A} \cdot \mathbf{s} + \mathbf{e} \in R_q^k.
    \]
    \item Khóa công khai là $pk = (\mathbf{t}, \rho)$. Chỉ $\rho$ được lưu thay vì toàn bộ $\mathbf{A}$ giúp tiết kiệm không gian đáng kể.
    \item Khóa bí mật là $sk = \mathbf{s}$ (trên thực tế, phiên bản IND-CCA2 gắn thêm $pk$, hash của $pk$ và một giá trị ngẫu nhiên $z$ vào $sk$ để phục vụ biến đổi Fujisaki-Okamoto trình bày ở mục sau).
\end{enumerate}

Về mặt toán học, $(\mathbf{A}, \mathbf{t})$ chính là một thể hiện của bài toán Module-LWE: từ $\mathbf{A}$ ngẫu nhiên và $\mathbf{t} = \mathbf{A} \cdot \mathbf{s} + \mathbf{e}$, việc tìm $\mathbf{s}$ tương đương giải Module-LWE, được giả định là khó ngay cả với máy tính lượng tử.

\subsection{Encapsulate}

Thuật toán ML-KEM.Encaps là thao tác phía người gửi. Đầu vào là khóa công khai $pk = (\mathbf{t}, \rho)$ của người nhận. Đầu ra gồm ciphertext KEM $c$ để gửi qua server và shared secret $K$ dài 256 bit để đưa vào quá trình dẫn xuất khóa. 

Các bước cốt lõi gồm:

\begin{enumerate}
    \item Sinh thông điệp ngẫu nhiên $m \in \{0,1\}^{256}$.
    \item Từ $\rho$ dựng lại ma trận $\mathbf{A}$ giống hệt bên sinh khóa.
    \item Lấy mẫu ba đại lượng ngẫu nhiên phụ thuộc $m$ và $pk$:
    \begin{itemize}
        \item Vector $\mathbf{r} \in R_q^k$ với các hệ số theo $B_{\eta_1}$.
        \item Vector nhiễu $\mathbf{e}_1 \in R_q^k$ theo $B_{\eta_2}$.
        \item Nhiễu scalar $e_2 \in R_q$ theo $B_{\eta_2}$.
    \end{itemize}
    Việc lấy mẫu là xác định theo $m$ thông qua hàm băm, nên cùng một $m$ luôn sinh ra cùng $(\mathbf{r}, \mathbf{e}_1, e_2)$. Đây là chi tiết quan trọng cho biến đổi Fujisaki-Okamoto.
    \item Tính hai thành phần của ciphertext:
    \[
    \mathbf{u} = \mathbf{A}^{\top} \cdot \mathbf{r} + \mathbf{e}_1 \in R_q^k,
    \]
    \[
    v = \mathbf{t}^{\top} \cdot \mathbf{r} + e_2 + \left\lceil \tfrac{q}{2} \right\rfloor \cdot m \in R_q,
    \]
    trong đó $\lceil q/2 \rfloor$ là phần nguyên gần nhất của $q/2$ (với $q = 3329$ thì $\lceil q/2 \rfloor = 1665$), và $\lceil q/2 \rfloor \cdot m$ có nghĩa là ``gắn'' thông điệp $m$ vào hệ số của đa thức bằng cách ánh xạ bit $0$ thành hệ số $0$ và bit $1$ thành hệ số $1665$.
    \item Nén ciphertext: cả $\mathbf{u}$ và $v$ đều được nén bằng cách cắt bớt các bit thấp nhất trong mỗi hệ số, giảm kích thước truyền đi. Ký hiệu $c = (\text{Compress}(\mathbf{u}), \text{Compress}(v))$.
    \item Dẫn xuất shared secret: $K = \text{KDF}(m \,\|\, \text{Hash}(c))$, với KDF là SHAKE-256. Đây chính là khóa cuối cùng mà AES-GCM sẽ dùng.
\end{enumerate}

Trực giác toán học là như sau. Nếu bỏ qua nhiễu, ta có $\mathbf{t} = \mathbf{A} \cdot \mathbf{s}$ và $v \approx \mathbf{t}^{\top} \mathbf{r} + \lceil q/2 \rfloor m = \mathbf{s}^{\top} \mathbf{A}^{\top} \mathbf{r} + \lceil q/2 \rfloor m$. Phía giải mã biết $\mathbf{s}$, có thể tính $\mathbf{s}^{\top} \mathbf{u} \approx \mathbf{s}^{\top} \mathbf{A}^{\top} \mathbf{r}$, rồi lấy $v - \mathbf{s}^{\top} \mathbf{u}$ sẽ gần bằng $\lceil q/2 \rfloor m$. Các nhiễu nhỏ bị tích lũy trong quá trình nhưng vẫn đủ nhỏ để phục hồi được $m$ bằng cách làm tròn.

\subsection{Decapsulate}

ML-KEM.Decaps là thao tác phía người nhận. Đầu vào là khóa bí mật $sk$ đang lưu cục bộ và ciphertext KEM $c = (\mathbf{u}, v)$ nhận từ người gửi. Đầu ra là shared secret $K$ giống với phía Encaps nếu ciphertext hợp lệ. Nếu ciphertext bị sửa đổi hoặc không hợp lệ, triển khai chuẩn không trả lỗi trực tiếp mà tạo một secret giả ngẫu nhiên, giúp tránh để lộ thông tin qua phản hồi giải mã.

Các bước gồm:

\begin{enumerate}
    \item Giải nén $\mathbf{u}$ và $v$ về dạng đa thức.
    \item Tính
    \[
    w = v - \mathbf{s}^{\top} \cdot \mathbf{u} \in R_q.
    \]
    \item Khôi phục thông điệp $m'$ bằng cách làm tròn từng hệ số của $w$: hệ số gần $0$ cho bit $0$, hệ số gần $\lceil q/2 \rfloor$ cho bit $1$. Cụ thể, với mỗi hệ số $w_i \in \{0, \ldots, q-1\}$ ta đặt
    \[
    m'_i = \left\lfloor \frac{2 \cdot w_i}{q} \right\rceil \bmod 2.
    \]
    \item Chạy lại toàn bộ Encaps với $m'$ để thu được $c'$ và so sánh với ciphertext đầu vào. Nếu $c' = c$, trả về
    \[
    K = \text{KDF}(m' \,\|\, \text{Hash}(c)).
    \]
    Nếu không, trả về
    \[
    K = \text{KDF}(z \,\|\, \text{Hash}(c)),
    \]
    với $z$ là giá trị ngẫu nhiên bí mật được gắn vào $sk$ lúc sinh khóa. Kiểm tra so sánh này chính là cốt lõi của biến đổi Fujisaki-Okamoto, bảo đảm IND-CCA2.
\end{enumerate}

\subsection{Phân tích tính đúng đắn}

Để thấy vì sao quá trình giải mã khôi phục lại được $m$, xét phép tính ở bước 2 của Decaps khi bỏ qua phần nén để đơn giản:
\[
w = v - \mathbf{s}^{\top} \mathbf{u}
  = \left( \mathbf{t}^{\top} \mathbf{r} + e_2 + \lceil q/2 \rfloor m \right) - \mathbf{s}^{\top} \left( \mathbf{A}^{\top} \mathbf{r} + \mathbf{e}_1 \right).
\]
Thay $\mathbf{t}^{\top} = (\mathbf{A} \mathbf{s} + \mathbf{e})^{\top} = \mathbf{s}^{\top} \mathbf{A}^{\top} + \mathbf{e}^{\top}$ ta được:
\[
w = \mathbf{s}^{\top} \mathbf{A}^{\top} \mathbf{r} + \mathbf{e}^{\top} \mathbf{r} + e_2 + \lceil q/2 \rfloor m - \mathbf{s}^{\top} \mathbf{A}^{\top} \mathbf{r} - \mathbf{s}^{\top} \mathbf{e}_1.
\]
Rút gọn thành
\[
w = \lceil q/2 \rfloor \cdot m + \underbrace{\mathbf{e}^{\top} \mathbf{r} - \mathbf{s}^{\top} \mathbf{e}_1 + e_2}_{\text{sai số tổng hợp } \delta}.
\]
Như vậy $w$ bằng $\lceil q/2 \rfloor m$ cộng một sai số $\delta$. Do $\mathbf{e}, \mathbf{r}, \mathbf{s}, \mathbf{e}_1, e_2$ đều có hệ số nhỏ (theo phân phối $B_\eta$ với $\eta \leq 2$), kỳ vọng và phương sai của từng hệ số của $\delta$ đều được giới hạn. Các tham số $n$, $q$, $\eta_1$, $\eta_2$ và sơ đồ nén được chọn sao cho với xác suất rất cao (cụ thể, $\Pr[\text{giải mã sai}] < 2^{-164}$ với ML-KEM-768), mỗi hệ số của $\delta$ có độ lớn nhỏ hơn $q/4$. Khi đó, phép làm tròn ở bước 3 của Decaps khôi phục chính xác từng bit của $m$:
\[
m'_i = \left\lfloor \frac{2 w_i}{q} \right\rceil = \left\lfloor \frac{2 (\lceil q/2 \rfloor m_i + \delta_i)}{q} \right\rceil = m_i \bmod 2.
\]
Đây chính là lý do ML-KEM vừa đảm bảo tính đúng đắn gần như hoàn hảo, vừa giữ được mức nhiễu đủ cho Module-LWE trở thành khó.

\subsection{Biến đổi Fujisaki-Okamoto và bảo mật IND-CCA2}

ML-KEM áp dụng biến đổi Fujisaki-Okamoto để nâng lên mức bảo mật IND-CCA2 - an toàn ngay cả khi kẻ tấn công có quyền truy cập tới một oracle giải mã.

Ý tưởng của biến đổi này là thay vì mã hóa ngẫu nhiên, các giá trị $(\mathbf{r}, \mathbf{e}_1, e_2)$ được dẫn xuất xác định từ chính thông điệp $m$ và khóa công khai $pk$ thông qua hàm băm:
\[
(\mathbf{r}, \mathbf{e}_1, e_2) = G(m, pk),
\]
với $G$ là hàm băm SHAKE-256. Nhờ đó, người giải mã sau khi khôi phục $m'$ có thể tự tính lại ciphertext $c'$ một cách xác định và so sánh với $c$. Nếu $c' \neq c$, kết luận rằng ciphertext đã bị sửa đổi và trả về một khóa giả ngẫu nhiên (được dẫn xuất từ giá trị bí mật $z$) thay vì báo lỗi. Cơ chế ``giả ngẫu nhiên thay vì báo lỗi'' này được gọi là implicit rejection, giúp tránh tấn công kênh kề dạng time-based khi kẻ tấn công cố gắng phân biệt giữa ciphertext hợp lệ và không hợp lệ.

Kết quả là bộ ba $(\text{KeyGen}, \text{Encaps}, \text{Decaps})$ đầy đủ của ML-KEM đạt mức an toàn IND-CCA2 trong Random Oracle Model lượng tử, được chứng minh reduction tới độ khó Module-LWE. Đây chính là mức bảo mật mạnh nhất cần thiết cho một KEM được sử dụng thực tế, trong đó kẻ tấn công có thể vừa có khóa công khai, vừa chủ động gửi ciphertext bất kỳ để quan sát phản hồi.

\subsection{Vì sao ML-KEM phù hợp với hệ thống nhắn tin}

Đối với một ứng dụng nhắn tin, ML-KEM có ba điểm mạnh chính. Thứ nhất, thuật toán đã được NIST chuẩn hóa trong FIPS 203, nên phù hợp hơn cho một hệ thống nghiêm túc so với việc tự chọn một đề xuất chưa ổn định. Thứ hai, đầu ra của ML-KEM là một shared secret ngắn, dễ đưa vào HKDF để kết hợp với X25519, Double Ratchet hoặc cơ chế MLS group. Thứ ba, thao tác encapsulate và decapsulate đủ nhanh để chạy trên client desktop mà không tạo độ trễ rõ rệt cho người dùng.

ML-KEM không dùng để mã hóa trực tiếp nội dung tin nhắn. Ciphertext KEM chỉ có nhiệm vụ vận chuyển một secret dùng cho bắt tay. Sau khi secret được tạo, nội dung tin nhắn vẫn nên được mã hóa bằng thuật toán đối xứng như AES-GCM. Cách tách vai trò này giúp hệ thống đơn giản hơn: KEM giải quyết bài toán thiết lập secret ban đầu, còn AEAD giải quyết bài toán mã hóa và xác thực từng message.

\subsection{Các biến thể và mức an toàn}

ML-KEM được cung cấp ở ba biến thể tương ứng với ba mức an toàn khác nhau. ML-KEM-512 với $k=2$ cung cấp mức an toàn tương đương AES-128, phù hợp cho các ứng dụng hạn chế về tài nguyên. ML-KEM-768 với $k=3$ cung cấp mức an toàn tương đương AES-192, là lựa chọn cân bằng giữa an toàn và hiệu năng, và được khuyến nghị cho đa số ứng dụng thực tế. ML-KEM-1024 với $k=4$ đạt mức an toàn tương đương AES-256 và phù hợp cho các ứng dụng có yêu cầu bảo mật cao nhất.

Về kích thước, ML-KEM-768 có khóa công khai 1184 byte, khóa bí mật 2400 byte, ciphertext 1088 byte và shared secret 32 byte. Các giá trị này lớn hơn đáng kể so với ECDH trên đường cong X25519 với khóa chỉ 32 byte, nhưng vẫn trong giới hạn có thể chấp nhận đối với một ứng dụng desktop hiện đại. Bù lại, thời gian thực thi của ML-KEM ở các thao tác keygen, encapsulate và decapsulate đều rất nhanh, thường chỉ trong khoảng vài trăm micro giây trên một CPU phổ thông.

\section{Forward secrecy, Double Ratchet và MLS}

\subsection{Forward secrecy và post-compromise security}

Forward secrecy là thuộc tính bảo đảm rằng việc lộ một khóa dài hạn ở thời điểm sau không tự động làm lộ các tin nhắn đã gửi trước đó. Trong nhắn tin, yêu cầu này quan trọng vì transcript có thể bị lưu lại rất lâu bởi mạng, server hoặc bản sao lưu. Nếu hệ thống dùng một khóa hội thoại cố định cho toàn bộ lịch sử, chỉ cần khóa đó bị lộ thì toàn bộ transcript cũ đều nguy hiểm.

Post-compromise security là thuộc tính mạnh hơn ở chiều ngược lại. Nếu attacker từng lấy được một phần trạng thái tạm thời trên thiết bị nhưng sau đó không còn kiểm soát thiết bị nữa, giao thức cần có khả năng tự phục hồi sau khi hai bên tiếp tục trao đổi secret mới. Hai thuộc tính này không đến từ một thuật toán đơn lẻ, mà đến từ cách tổ chức vòng đời khóa: khóa message chỉ dùng một lần, chain key được cập nhật liên tục, và các secret mới được đưa vào state theo thời gian.

\subsection{Double Ratchet trong hội thoại hai người}\label{subsec:theory-dr}

Double Ratchet là cơ chế đạt forward secrecy ở mức từng tin nhắn cho hội thoại hai người, được phổ biến qua giao thức Signal. Sau bước bắt tay ban đầu, hai phía cùng có một \emph{root key} bí mật chung. Toàn bộ phần còn lại của giao thức xoay quanh việc liên tục biến đổi root key này thành các khóa con dùng một lần, theo hai ``bánh cóc'' (ratchet) lồng nhau: một bánh cóc đối xứng chạy theo từng tin, và một bánh cóc Diffie-Hellman chạy theo từng lượt đổi chiều. Sở dĩ gọi là bánh cóc vì giống cơ cấu cơ khí chỉ quay được một chiều: trạng thái khóa chỉ tiến tới, không thể quay ngược lại để khôi phục các giá trị cũ.

\textbf{Trạng thái phiên.} Ở mỗi phía, mỗi cuộc hội thoại được mô tả bằng một trạng thái phiên  ($\mathsf{st}$) - là toàn bộ vật liệu khóa và bộ đếm cần để mã hóa tin kế tiếp và giải mã tin nhận về:
\[
  \mathsf{st} = \bigl(RK,\ CK_s,\ CK_r,\ (pk_{\text{self}}, sk_{\text{self}}),\ pk_{\text{remote}},\ N_s,\ N_r,\ PN,\ \mathsf{MK}_{\text{skip}}\bigr).
\]
Ý nghĩa và vai trò của từng thành phần:
\begin{itemize}
  \item $RK$ - \emph{root key}, bí mật gốc của phiên (ban đầu do bắt tay PQXDH gieo). Nó chỉ thay đổi khi có bánh cóc DH nên đóng vai trò trục ổn định giữa các lượt đổi chiều, và là nguồn dẫn ra mọi chain key.
  \item $CK_s,\ CK_r$ - hai \emph{chain key} cho chiều gửi và chiều nhận, mỗi tin tiến một bước qua $\mathrm{KDF\_CK}$. Vì hai phía cùng một $RK$ nhưng tráo vai nên chain gửi của bên này đúng bằng chain nhận của bên kia.
  \item $(pk_{\text{self}}, sk_{\text{self}})$ - cặp khóa X25519 ratchet hiện hành của mình; khóa công khai $pk_{\text{self}}$ được đính vào header mỗi tin gửi đi.
  \item $pk_{\text{remote}}$ - khóa ratchet công khai mới nhất của đối phương. Khi header một tin đến mang khóa khác giá trị này, đó là tín hiệu để thực hiện một bước DH ratchet.
  \item $N_s,\ N_r$ - số tin đã gửi trong chain gửi hiện tại, và số thứ tự tin mong đợi tiếp theo ở chain nhận.
  \item $PN$ - (\emph{previous N}) độ dài chain gửi trước lần xoay DH ratchet gần nhất. Nhờ $PN$ đặt trong header, phía nhận biết phải sinh và lưu nốt bao nhiêu khóa của chain cũ trước khi chuyển sang chain mới, để các tin cũ đến muộn vẫn giải mã được.
  \item $\mathsf{MK}_{\text{skip}}$ - bảng message key của các tin bị bỏ qua (đến muộn hoặc lệch thứ tự), có giới hạn kích thước để chống bị làm cạn bộ nhớ.
\end{itemize}
Mọi khóa trong $\mathsf{st}$ đều dài 32 byte (root key, chain key, message key và khóa X25519). Cài đặt thực tế còn bổ sung vài trường để ràng buộc metadata (khóa mã hóa header, định danh phiên và hội thoại, số thứ tự định danh) cùng cơ chế chống nhận lặp, nhưng không làm thay đổi bản chất hai bánh cóc.

\textbf{Bánh cóc đối xứng (symmetric-key ratchet).} Từ root key, mỗi phía dẫn xuất hai \emph{chain key}: một cho chiều gửi và một cho chiều nhận. Mỗi khi cần xử lý một tin, chain key được đưa qua một hàm dẫn xuất một chiều $\mathrm{KDF\_CK}$ để tách ra đồng thời một \emph{message key} dùng một lần và một chain key kế tiếp:
\[
  MK_i = \mathrm{HMAC}(CK_i, \texttt{0x01}), \qquad
  CK_{i+1} = \mathrm{HMAC}(CK_i, \texttt{0x02}).
\]
Message key $MK_i$ chỉ dùng để mã hóa hoặc giải mã đúng một tin rồi bị xóa; chain key tiến từ $CK_i$ sang $CK_{i+1}$ và giá trị cũ cũng bị xóa ngay. Vì $\mathrm{HMAC}$ là một chiều, biết $CK_{i+1}$ hay $MK_i$ không cho phép suy ngược ra $CK_i$ hoặc bất kỳ message key nào trước đó. Đây chính là nguồn gốc của forward secrecy ở mức tin nhắn: việc lộ trạng thái hiện tại không làm lộ các tin đã gửi trước.

\textbf{Bánh cóc Diffie-Hellman (DH ratchet).} Riêng bánh cóc đối xứng thì không tự phục hồi: nếu attacker lấy được chain key hiện tại, hắn có thể tính mọi message key về sau trong cùng một chiều. Để vá điểm này, mỗi phía giữ thêm một cặp khóa DH tạm thời và đính kèm khóa công khai DH hiện hành của mình vào header mỗi tin. Khi một phía nhận được một khóa DH công khai mới của đối phương, nó thực hiện một bước DH ratchet: tính một secret DH mới rồi trộn vào root key để sinh ra root key và chain key mới qua hàm $\mathrm{KDF\_RK}$:
\[
  RK' \,\Vert\, CK' = \mathrm{HKDF}\bigl(\,\text{salt}=RK,\ \ \text{IKM}=\mathrm{DH}(sk_{\text{self}}, pk_{\text{remote}}),\ \ \texttt{info},\ 64\,\bigr).
\]
Đầu ra 64 byte được cắt làm đôi: 32 byte đầu là root key mới $RK'$, 32 byte sau là chain key mới $CK'$ cho chiều vừa thiết lập; \texttt{info} là một nhãn cố định gắn với ứng dụng. Bên trong, $\mathrm{HKDF}$ gồm bước Extract (dùng $RK$ làm salt để trộn DH secret thành một khóa giả ngẫu nhiên) rồi bước Expand (kéo dài thành 64 byte có cấu trúc). Sau đó phía này tự sinh một cặp khóa DH mới cho riêng mình và gửi kèm khóa công khai trong tin tiếp theo, để đối phương cũng thực hiện bước ratchet tương ứng. Nhờ liên tục bơm secret DH tươi vào root key, giao thức đạt post-compromise security: chỉ cần hai bên còn trao đổi được một lượt khóa DH mới mà attacker không quan sát, trạng thái sẽ ``lành lại'' và các tin sau đó an toàn trở lại.

\textbf{Thuật toán gửi và nhận một tin.} Khi gửi, phía gửi tiến chain key gửi để lấy message key, gắn header rồi mã hóa bằng AEAD với header làm dữ liệu xác thực kèm theo (AAD):
\begin{enumerate}
  \item $(CK_s,\ MK) \leftarrow \mathrm{KDF\_CK}(CK_s)$ và tăng $N_s$;
  \item dựng header $H = (pk_{\text{self}},\ PN,\ N_s)$;
  \item $c = \mathrm{AEAD.Enc}\bigl(MK;\ \text{plaintext};\ \mathsf{AAD}=H\bigr)$, rồi xóa ngay $MK$.
\end{enumerate}
Khi nhận một tin có header $H' = (pk',\ PN',\ N')$: nếu $pk'$ khác $pk_{\text{remote}}$ đang lưu thì trước hết sinh và cất các message key còn thiếu của chain cũ tới chỉ số $PN'$, rồi thực hiện một \emph{bước DH ratchet} (tính $RK,\ CK_r$ mới theo công thức trên, sinh cặp khóa DH mới và đặt lại $N_s=N_r=0$). Tiếp đó phía nhận tiến chain nhận tới chỉ số $N'$, lưu các message key trung gian vào $\mathsf{MK}_{\text{skip}}$, lấy ra $MK$ rồi giải mã và kiểm tra tag AEAD trên cả header lẫn ciphertext. Việc đưa header vào AAD ràng buộc mỗi ciphertext với đúng khóa DH và chỉ số của nó, nên kẻ tấn công không thể tráo header để đánh lừa bước ratchet.

\textbf{Tin lệch thứ tự và chống lặp.} Vì mạng có thể giao tin sai thứ tự, header còn mang hai bộ đếm: số thứ tự của tin trong chain hiện tại và số tin của chain trước đó. Khi nhận một tin có số thứ tự vượt trước, phía nhận lần lượt sinh và lưu tạm các message key bị bỏ qua (\emph{skipped keys}) trong một giới hạn an toàn để vẫn giải mã được tin đến muộn, đồng thời từ chối ciphertext đã dùng nhằm chống replay. Trong đồ án, giới hạn này là 1000 khóa để cân bằng giữa khả năng chịu lệch thứ tự và rủi ro bị làm cạn bộ nhớ.

\textbf{Liên hệ hậu lượng tử.} Bản thân bánh cóc dùng X25519 nên không kháng lượng tử; tính kháng lượng tử của hội thoại nằm ở bước bắt tay ban đầu, nơi root key đầu tiên được gieo từ một bắt tay kiểu PQXDH có trộn thêm shared secret ML-KEM (trình bày ở chương thiết kế). Nói cách khác, ML-KEM bảo vệ bí mật của điểm khởi đầu, còn Double Ratchet bảo vệ tiến trình của hội thoại theo thời gian.

\subsection{Messaging Layer Security cho hội thoại nhóm}\label{subsec:theory-mls}

Nếu áp dụng Double Ratchet riêng cho từng cặp thành viên trong một nhóm $n$ người, mỗi tin phải mã hóa lại $n-1$ lần và số phiên cần đồng bộ tăng theo $O(n^2)$; việc thêm/xóa thành viên và giữ đúng các thuộc tính an toàn trở nên rất khó. Messaging Layer Security (MLS, RFC 9420) giải bài toán này bằng một \emph{group state} chung, cho phép làm mới khóa nhóm với chi phí chỉ $O(\log n)$ cho mỗi thao tác.

\textbf{Cây ratchet và TreeKEM.} Trái tim của MLS là một cây nhị phân cân bằng gọi là \emph{ratchet tree}: mỗi thành viên là một lá, mỗi nút (kể cả nút trong) gắn với một cặp khóa bất đối xứng. Bất biến cốt lõi của TreeKEM là: một thành viên biết khóa bí mật của đúng những nút nằm trên đường đi trực tiếp từ lá của mình lên gốc (\emph{direct path}), và biết khóa công khai của mọi nút còn lại. Khóa bí mật của nút gốc do đó là một secret mà toàn bộ và chỉ các thành viên hiện tại mới cùng suy ra được.

Cây được đánh số theo thứ tự duyệt giữa (in-order): với $n$ lá, cây có $2n-1$ nút, lá ở các chỉ số chẵn còn nút trong ở các chỉ số lẻ; cha và hai con của một nút được xác định thuần túy bằng số học trên chỉ số, nên mọi thành viên cùng nhìn thấy một cấu trúc cây thống nhất mà không phải trao đổi thêm. Direct path của một lá là dãy nút từ cha trực tiếp lên gốc; \emph{copath} là dãy các nút anh em của những nút đó (phần cây con ở phía đối diện).

\textbf{Cập nhật đường đi (path update).} Khi một thành viên làm mới khóa của mình (thao tác Update) hoặc khi có thay đổi thành viên, nó sinh một secret lá ngẫu nhiên $s_0$ rồi dẫn ra một chuỗi \emph{path secret} dọc theo direct path; tại mỗi mức, path secret vừa sinh ra cặp khóa HPKE của nút đó, vừa dẫn tiếp lên mức cha:
\[
  s_{i} = \mathrm{DeriveSecret}(s_{i-1},\ \mathtt{path}), \qquad
  (pk_i, sk_i) = \mathrm{DeriveKeyPair}\bigl(\mathrm{DeriveSecret}(s_i,\ \mathtt{node})\bigr),
\]
trong đó $\mathrm{DeriveSecret}(s, L) = \mathrm{ExpandWithLabel}(s, L, \varepsilon, N_h)$ là HKDF-Expand có gắn nhãn $L$ (độ dài đầu ra $N_h$ byte). Để các thành viên khác cập nhật theo, với mỗi nút $v$ trên direct path, người cập nhật mã hóa path secret tương ứng cho \emph{tập độ phân giải} (resolution) của copath bằng HPKE:
\[
  \mathsf{ct}_v = \mathrm{HPKE.Enc}\bigl(pk_{\mathrm{res}(v)},\ s_{i(v)}\bigr).
\]

Một thành viên khác giải mã được $\mathsf{ct}_v$ tại đúng điểm hai đường đi gặp nhau (nút tổ tiên chung thấp nhất), rồi tự dẫn tiếp các path secret còn lại lên gốc bằng chính công thức trên. Nhờ chỉ phải mã hóa cho $O(\log n)$ nút thay vì cho từng người, chi phí một lần cập nhật giảm từ tuyến tính xuống còn logarit theo số thành viên.

\textbf{Lịch khóa và epoch.} Mỗi cấu hình nhóm là một \emph{epoch}. Sau một thao tác, secret rút ra từ gốc cây (\emph{commit secret}) được trộn với init secret của epoch trước qua hai bước HKDF để tạo joiner secret rồi epoch secret:
\[
  \text{joiner}_n = \mathrm{KDF.Extract}\bigl(\text{init}_{n-1},\ \text{commit}_n\bigr),
\]
\[
  \text{epoch}_n = \mathrm{ExpandWithLabel}\bigl(\mathrm{KDF.Extract}(\text{joiner}_n,\ \text{psk}),\ \mathtt{epoch},\ \mathrm{GroupContext}_n,\ N_h\bigr),
\]
trong đó $\text{psk}$ là pre-shared key (bằng 0 nếu không dùng) và $\mathrm{GroupContext}_n$ buộc khóa vào đúng cấu hình nhóm ở epoch $n$. Từ $\text{epoch}_n$, mỗi secret chuyên biệt được dẫn ra bằng $\mathrm{DeriveSecret}(\text{epoch}_n, L)$ với một nhãn $L$ riêng, trong đó đáng chú ý là:
\begin{itemize}
  \item \texttt{encryption} - gốc của \emph{secret tree} sinh khóa mã hóa tin ứng dụng;
  \item \texttt{sender data} - bảo vệ metadata định danh người gửi của mỗi tin;
  \item \texttt{confirm} và \texttt{membership} - khóa xác thực transcript và tư cách thành viên trong handshake;
  \item \texttt{init} - chính là $\text{init}_n$, được nối làm đầu vào cho key schedule của epoch kế tiếp.
\end{itemize}
Việc init secret của epoch này nối tiếp vào epoch sau tạo thành một bánh cóc bất đối xứng ở quy mô nhóm: mỗi lần đổi epoch, toàn bộ khóa cũ bị bỏ (forward secrecy) và entropy mới từ commit secret được bơm vào (post-compromise security).

\textbf{Cây secret và bánh cóc ứng dụng trong một epoch.} Riêng việc đổi epoch chỉ cho forward secrecy ở mức thô. Để mỗi tin trong cùng một epoch cũng có khóa riêng, secret nhãn \texttt{encryption} được dùng làm gốc của một cây nhị phân thứ hai gọi là \emph{secret tree}; mỗi lá ứng với một thành viên và cho ra một chuỗi ratchet đối xứng riêng. Tại thế hệ $j$, khóa và nonce AEAD cùng secret kế tiếp được dẫn ra bằng HKDF có nhãn:

\begin{itemize}
  \item Khóa AEAD: $k_j = \mathrm{ExpandWithLabel}(r_j, \mathtt{key}, j, N_k)$;
  \item Nonce AEAD: $n_j = \mathrm{ExpandWithLabel}(r_j, \mathtt{nonce}, j, N_n)$;
  \item Secret kế tiếp: $r_{j+1} = \mathrm{ExpandWithLabel}(r_j, \mathtt{secret}, j, N_h)$.
\end{itemize}

Mỗi tin dùng một cặp $(k_j, n_j)$ riêng rồi $r_j$ tiến tới $r_{j+1}$ và bị xóa, cho forward secrecy theo từng tin ngay trong một epoch. Như vậy MLS đạt cả hai thuộc tính theo hai trục bổ sung nhau: bánh cóc đối xứng trong secret tree lo forward secrecy giữa các tin, còn bánh cóc bất đối xứng qua các epoch (path update) lo post-compromise security ở mức nhóm.

\textbf{Đề xuất, Commit và Welcome.} Thay đổi nhóm được biểu diễn bằng các \emph{proposal}: Add (thêm thành viên qua KeyPackage của họ), Remove (xóa thành viên), Update (làm mới khóa của một thành viên). Một \emph{Commit} gói một tập proposal lại, thực hiện path update và đẩy nhóm sang epoch mới. Thành viên mới nhận một thông điệp \emph{Welcome} chứa đủ secret (mã hóa riêng cho KeyPackage của họ) để bước vào đúng epoch; thành viên bị Remove không nằm trong tập nhận secret của epoch kế tiếp nên lập tức mất khả năng giải mã.

\textbf{Tách public state và secret state.} MLS phân biệt rõ phần công khai và phần bí mật của group state. Các phần public như GroupInfo, cấu trúc ratchet tree (chỉ gồm khóa công khai), epoch và danh sách credential cho phép kiểm tra tính nối tiếp và hợp lệ của nhóm mà không lộ application secret; các secret như epoch secret và path secret chỉ tồn tại ở client. Nhờ vậy một server không tin cậy vẫn có thể lưu trữ, chuyển tiếp và kiểm tra ràng buộc membership dựa trên public state, nhưng không bao giờ giải mã được tin nhắn nhóm.

\section{Giao thức TLS và vai trò trong kiến trúc bảo mật}

Mặc dù mã hóa đầu cuối đã bảo vệ nội dung tin nhắn khỏi server, TLS ở tầng vận chuyển vẫn có ý nghĩa quan trọng trong một kiến trúc bảo mật nhiều lớp. TLS che nội dung gói JSON khỏi người nghe lén trên mạng và giúp xác thực máy chủ ở cấp độ hạ tầng. Tuy nhiên TLS không che hoàn toàn metadata như địa chỉ kết nối, thời điểm truyền và kích thước gói quan sát được trên mạng. Khi kết hợp mã hóa đầu cuối với TLS, hệ thống có hai lớp phòng vệ khác nhau thay vì chỉ phụ thuộc vào một lớp.

Trong đồ án, tầng TLS giữa client và server được thiết lập với chứng chỉ X.509. Toàn bộ các gói JSON được gửi trong kênh TLS này, bao gồm cả những gói đã chứa ciphertext đầu cuối. Nhờ đó, kẻ tấn công trên đường truyền không đọc được nội dung JSON. Dù vậy, server vẫn là điểm cuối của TLS, nên nội dung chat vẫn phải được mã hóa thêm ở tầng ứng dụng bằng E2EE.

\section{Các công nghệ triển khai}

\subsection{Ngôn ngữ Python và PyQt5}

Python là ngôn ngữ được lựa chọn để hiện thực hóa phía client nhờ vào sự kết hợp giữa cú pháp rõ ràng, hệ sinh thái thư viện phong phú và khả năng tích hợp tốt với các thành phần mật mã. Đối với phần giao diện, PyQt5 là framework trưởng thành và ổn định, cung cấp binding đầy đủ của thư viện Qt cho Python, cho phép xây dựng giao diện desktop native trên cả Windows, Linux và macOS từ cùng một mã nguồn.

PyQt5 hỗ trợ kiến trúc signal-slot, rất phù hợp cho các ứng dụng mạng bất đồng bộ, nơi các sự kiện từ luồng mạng cần được truyền về luồng giao diện một cách an toàn. Framework cũng cung cấp nhiều thành phần trực quan sẵn có như QScrollArea, QWidget, QListWidget, QFileDialog và nhiều widget khác để xây dựng message list, danh sách liên lạc, chọn file và các dialog chức năng.

\subsection{Ngôn ngữ C và thư viện hệ thống}

Phía máy chủ được triển khai bằng ngôn ngữ C để đảm bảo hiệu năng cao, mức sử dụng tài nguyên thấp và khả năng kiểm soát bộ nhớ chi tiết. Một máy chủ nhắn tin phải xử lý đồng thời nhiều kết nối mạng, mỗi kết nối có thể truyền tải các gói tin kích thước rất khác nhau từ tin nhắn văn bản ngắn vài byte cho đến file đa phương tiện nhiều megabyte. C cho phép quản lý bộ nhớ động một cách chính xác, tránh được chi phí của garbage collector và tận dụng tối đa các hệ thống lập lịch luồng của hệ điều hành.

Server trong đồ án sử dụng các thư viện hệ thống chuẩn của Linux bao gồm pthread cho mô hình đa luồng, socket BSD cho lập trình mạng, và OpenSSL cho tầng TLS. Ngoài ra, MySQL Connector/C được sử dụng để kết nối tới cơ sở dữ liệu, và một parser JSON nhẹ được viết thủ công để phân tích các gói tin mà không phụ thuộc vào thư viện bên ngoài nặng nề.

\subsection{Thư viện Open Quantum Safe}

Để tích hợp ML-KEM vào ứng dụng, đồ án sử dụng thư viện Open Quantum Safe, một dự án mã nguồn mở cung cấp hiện thực hóa các thuật toán mật mã hậu lượng tử đang được chuẩn hóa hoặc trong quá trình đánh giá. Thư viện lõi liboqs được viết bằng C với các tối ưu hóa cho nhiều kiến trúc CPU khác nhau, đồng thời cung cấp binding cho các ngôn ngữ phổ biến bao gồm Python, C++, Java và Go.

Trong đồ án, binding Python oqs-python được sử dụng ở phía client để thực hiện các thao tác KEM cần thiết như sinh khóa, encapsulate và decapsulate. Việc sử dụng một thư viện đã được kiểm định và duy trì tích cực là điều quan trọng, bởi mật mã hậu lượng tử có nhiều chi tiết cài đặt tinh tế mà việc tự viết lại từ đầu có thể dẫn đến các lỗi bảo mật khó phát hiện.

\subsection{Thư viện cryptography và các thao tác AES-GCM}

Đối với mã hóa đối xứng AES-GCM và hàm dẫn xuất khóa PBKDF2, thư viện cryptography của Python được lựa chọn. Đây là thư viện mật mã chính thức được cộng đồng Python Cryptographic Authority phát triển, cung cấp cả giao diện cấp cao dễ sử dụng và giao diện cấp thấp cho người cần kiểm soát chi tiết. Thư viện được xây dựng trên nền OpenSSL với các ràng buộc CFFI, đảm bảo hiệu năng gần bằng cài đặt C thuần và tự động hưởng lợi từ các tăng tốc phần cứng như AES-NI trên các CPU hỗ trợ.

\subsection{Cơ sở dữ liệu MySQL}

MySQL được chọn làm cơ sở dữ liệu quan hệ do tính phổ biến, tài liệu phong phú và hiệu năng tốt trong các tác vụ ghi-đọc nhiều. Ứng dụng nhắn tin có mô hình dữ liệu tương đối đơn giản với các thực thể như người dùng, bạn bè, hội thoại và tin nhắn, rất phù hợp với mô hình quan hệ. MySQL cũng hỗ trợ kiểu dữ liệu LONGTEXT lên tới 4 GB cho mỗi bản ghi, rất có ích khi lưu trữ tin nhắn chứa file đã được mã hóa và mã hóa base64 với độ dài có thể lên tới vài chục megabyte.

\subsection{Docker và Docker Compose}

Để đơn giản hóa việc triển khai và đảm bảo môi trường chạy nhất quán giữa các máy phát triển khác nhau, đồ án sử dụng Docker làm công cụ container hóa cùng Docker Compose để điều phối nhiều container. Server C và MySQL được đóng gói thành hai container riêng biệt với các volume lưu trữ bền vững cho dữ liệu cơ sở dữ liệu. Với một lệnh docker compose up duy nhất, toàn bộ môi trường máy chủ có thể được khởi động, tạo điều kiện thuận lợi cho việc phát triển, kiểm tra nội bộ và trình diễn.
